\section{Risks \& Mitigation Strategies}

There are several key risks associated with my proposed project. Here, I discuss how I plan to address them:

\begin{itemize}
   \item \textbf{Design of constraint embedding mechanism:} The representation of dynamic constraints as embeddings for diffusion conditioning is non-trivial. I will start with low-dispersion sampling of the robot state space with binary vectors indicating problem-specific constraint satisfaction---a generalization of PRESTO to accommodate dynamic constraints \cite{seo2024presto}. If this proves insufficient for complex constraint manifolds, I will explore progressively more sophisticated embeddings: 1) addressing the high-dimensionality of the state space by creating submaps as channels for a richer representation \cite{krizhevsky2012imagenet}; 2) employing PointNet-style encoders that are invariant to the permutation of constraint samples \cite{qi2017pointnet}; and, 3) investigating graph-based architectures to capture relationships between constraint-satisfying configurations \cite{hamilton2017inductive}.
   \item \textbf{Insufficient task complexity:} Initial test cases (payload transport in clutter-free environments) may not be complex enough to elicit the benefits of my proposed approach. To mitigate this risk, I will systematically increase task complexity by introducing collision objects, varying payload characteristics, and implementing more complex dynamic constraints (e.g., end-effector acceleration constraints for liquid handling \cite{ichnowski2022gomp,abderezaei2024clutter}.
   \item \textbf{Optimization problem design:} Designing optimization problems for payload and liquid handling tasks can be challenging, since it involves hyperparameter tuning, constraint linearization, and applying trust regions, which I only have a cursory knowledge of \cite{schulman2013finding}. I have already implemented have a slow version of optimization for payload transport which I will extend to accommodate liquid handling constraints \cite{pasricha2025payload}. I will also need to incorporate collision-checking, however there is extensive past work in this domain which will be useful \cite{ichnowski2020deep}.
\end{itemize}

